\section{Chaos in Mathematics}


\hfill

These are some reflections on what chaos means in the context of mathematics and physics. Chaos must be understood as
the absence, or privation, of order, not as the opposite of order in any active sense.

To explore what chaos may mean from the mathematical point of view, we can consider an irrational real number between
zero and one in its decimal expansion to try to represent a number representing true lack of order, that is, the
ultimate random number. A rational number is not random, since its digits will repeat after a certain point. The
ultimate random number will have two properties:

\begin{enumerate}
\item \textbf{Not Computable}: There is no algorithm to determine the sequence of digits 
\item \textbf{Normal}: Every digit, and sequence of digits, is evenly represented in the decimal expansion. 
\end{enumerate}
Numbers such as ${\surd}$2 and $\pi $ fail the first condition since there are mathematical formulas to compute the
decimal expansion to any desired precision.

In the decimal (or any base) expansion of normal numbers, each digit will appear with the same frequency, i.e., 10\%.
Furthermore, each sequence of digits will also appear with its corresponding frequency. That is, 10 through 99, 100
through 999, and so on, will also appear 1\% and 0.1\% of the time respectively. $\sqrt 2$ and $\pi $ are believed to
be normal, but there is no proof. The mathematician Gregory Chaitin defined a number $\Omega $,
Chaitin's constant, that has both properties.

Condition 1 is necessary because an algorithm is the hidden order behind a seemingly random sequence of numbers.
Condition 2 is necessary because complete chaos would have to include all possibilities.

In our random number, then, each digit is independent in the sense that its value does not depend on the digit, or
sequence of digits, that precede it. Investigating this number, we will observe at some point the sequence of numbers
from 1 to 1 million. This poses a dilemma. Do we conclude, then, that we have discovered a region of order within that
alleged chaos, or are we convinced that it is truly random? This is not merely speculative as it has real world
applications.

\begin{itemize}
\item In Intelligent Design theory, recognized patterns are presumed to result from a designer rather than random
variation. 
\item In the search for extra-terrestrial intelligence, cosmic radio waves are scanned to look for patterns hidden
within the white noise. 
\item In the multiverse theory of physics, all possible worlds exist and we just ``happen" to live in one of them that
exhibits the necessary degree of order to support human life. 
\item Then, there is the Infinite Monkey Theorem that claims that enough monkeys with typewriters, given enough time,
will produce the works of Shakespeare. 
\end{itemize}
We have previously asked this question about whether each moment in time is independent of any previous moment or if
there is a deep connection between the past and the present. The first is chaos; the second may be either order or
chaos deceptively ordered. There is no way, apparently, within a chaotic situation to determine if order is accidental
or real.

Two of the conditions of manifestation are form and matter. Chaos can only be associated with prime matter, Prakriti,
without property and infinitely receptive to any form. Matter, in this sense, can never ``exist", that is, it can never
be a possibility of manifestation. Hence, Chaitin's constant, is only a thought construction and
not a reality. That is why Guenon says that at the heart of the Kali Yuga, where chaos reigns, the state of total
disorder can never be reached.

Thus, form is a necessary condition of manifestation; it is transcendent to matter and we know it through our will and
intelligence. I design, hence I recognize design; I communicate, hence I recognize a communication; I create, hence I
recognize an artistic creation.



\flrightit{Posted on 2012-10-04 by Cologero }

\begin{center}* * *\end{center}

\begin{footnotesize}\begin{sffamily}



\texttt{Master Dogen on 2012-10-05 at 02:59 said: }

``That is why Guenon says that at the heart of the Kali Yuga, where chaos reigns, the state of total disorder can never
be reached."

This is a heartening thing to ponder.


\hfill

\texttt{william zeitler on 2012-10-05 at 09:52 said: }

``Chaos must be understood as the absence, or privation, of order, not as the opposite of order in any active sense."
I'm not sure that's necessarily true. Much of mathematical chaos theory is
about observing that there is an `order' that we simply
don't understand (e.g. fractals and `strange
attractors'). Also, related to your article might be Godel's ``Incompleteness
Theorem"–that for any non-trivial formal system there will always be statements that can neither be proven true nor
false. In other words, outside the light from the streetlight of, say, geometry, there will always be a geometrical
darkness we can't penetrate. (At least without cognitively stepping outside/above the geometrical
system itself.)


\hfill

\texttt{Cologero on 2012-10-05 at 11:28 said: }

Don't be misled by verbal similarity. I defined precisely what I meant by true chaos, in which
there is no possibility of a hidden order, except incidentally. But I agree with you in this sense that a real disorder
is not manifestable as ``there is an `order' that we simply
don't understand".


\hfill

\texttt{Andrew on 2012-10-05 at 18:21 said: }

It's not really heartening, as Prakriti is constituent to manifestation but requires Purusha to
take on any qualities. Knowing the impossible as impossible ought to be emotionally neutral.


\hfill

\texttt{Andrew on 2012-10-08 at 03:47 said: }

Bhakti Anand Goswami mentions Meditations on the Tarot in his video:

Tarot \& Ancient Symbols of God's Personal Form – Bhakti Ananda Goswami

\url{http://www.youtube.com/watch?v=-TBUn1Tog20\&feature=plcp}

He mentions the churning of the milk ocean as the ``perturbation of prakiti".


\hfill

\texttt{Master Dogen on 2012-10-08 at 06:03 said: }

It was a colloquialism, reflecting a feeling I sometimes have, of exasperation and astonishment at the ugliness and,
well, chaos of the modern world. When I have the wrong-colored glasses on, I sometimes wonder if literally everything
is going to fall apart. Perhaps it's not a very insightful comment, but even those of us seeking
the original Tradition exist in specific bodies with specific tendencies and paths, that result in specific emotional
reactions at different stages of those paths. 

If one is not seeing things clearly, then learning something that is indeed truthful can come as a great relief
— thus heartening. Have you never traveled from error into truth and enjoyed the process? If what
Cologero points out in this little article is SO self-evident that it has literally zero power to shed light for any of
us, then I suppose it's time for Cologero to stop writing them. However, at least for me, that is
not the case.

Furthermore, I can't claim that I KNOW the impossible as impossible. I read about Purusha and
Prakriti; I meditate on it; but it's not like I have succeeded in slaying the ego and manifesting
the eternal in my moment to moment existence. I don't suppose I would have much use for reading
Gornahoor if I were already totally secure in all my understandings of Tradition.


\end{sffamily}\end{footnotesize}
